%
%      Brno University of Technology
%      Faculty of Information Technology
%
%      MSc Thesis	2009/2010
%
%      Design of S-Boxes Using Genetic Algorithms
%
%      Bedrich Hovorka
% -----------------------------------------------------
%
In the design phase, a combination of C++ and Python languages was selected. Subsequently, in C++ with the help of
the GAlib library, the more time-consuming operations were implemented,
namely the evolution process, S-box representations,
and calculations of security criteria from S-box outputs.
For evolutionary experiments, algorithms that are not
part of the GAlib library were also chosen. Furthermore, an interface between
Python and C++ was implemented so that experiments could be launched from Python and the
resulting statistics could also be processed there.

The classes in C++ are divided into those that concern S-box representations and those that
perform searching. Any S-box representation can be used in any
search algorithm except the EDA algorithm, which by definition works with
numerical alleles; in this implementation only with binary ones.
In the following text, interesting components of the implementation are selected.

\section{Implementation of Evolutionary Algorithms}
First, the missing search algorithms were implemented. All
search algorithms are inherited from the base class \sourceRef{SboxSearchAlgorithmBase},
which inherits from \sourceRef{GASimpleGA} from GAlib. In this way, it was possible to use a more general
interface for search algorithms. Most of the logic of the following algorithms was
implemented in the overridden method \sourceRef{step()},
which performs one generation of the evolutionary process, from the superclass
\sourceRef{GASimpleGA} in GAlib. The following text describes the implemented
methods and auxiliary objects that are used during the execution of one generation
specifically for a particular algorithm.

\subsubsection{EDA Algorithms}
The EDA algorithm class extends the base class with a model property; it can be of two
types given by the object class. The base class \sourceRef{EdaModel} implements the statistical
UMDA model, while its subclass \sourceRef{BmdaModel} also uses simple allele occurrence probabilities
for its calculations. In the class \sourceRef{EstimationOfDistributionAlgorithm}
the general logic of the EDA algorithm is defined, and the responsibility for
creating the statistical model and generating new individuals is transferred
to an instance of the class \sourceRef{EdaModel}.

In the class \sourceRef{EdaModel}, two main polymorphic methods were implemented,
named
\sourceRef{learnStructure}, which simply iterates through all individuals and records the counts
of one-bits in an array of the length of the binary chromosome,
and \sourceRef{sampleModel}, which uses this array to generate individuals
with probability $P(genome[i] = 1) = count[i]/popSize$.

The class \sourceRef{BmdaModel} uses the class \sourceRef{ContingentTable} for creating the statistical model
and a forest\footnote{a set of trees}
created using instances of the \sourceRef{BmdaNode} classes for its storage. In the method \sourceRef{learnStructure}, it first calculates
contingency tables of dependencies between loci. Contingency tables for which
the hypothesis "dependent" of the $\chi^2$-test holds are sorted according to the value
of this $\chi^2$-test. If all loci are independent, the algorithm
terminates with only the calculated UMDA model. Otherwise, dependency trees
are generated from the sorted contingency tables\footnote{The algorithm for creating this forest is given e.g. in \cite{evo}}.
In the method \sourceRef{sampleModel}, either generation of individuals in the superclass is selected according to the type of the currently stored model,
or when generating each individual,
the forest is traversed and values at individual loci are calculated based on
the conditional probabilities encountered in nodes on tree branches.

\bigskip

The superclass of multi-criteria
algorithms\footnote{in the source code \sourceRef{MulticriterialGeneticAlgorithm}}
manages the set of required criteria. It also contains routines that
both multi-criteria algorithms use. These include, for example,
decision about dominance and calculation of the non-dominated elite using the
continuously updated approach method.

\subsubsection{VEGA Algorithm}
In the implementation of the VEGA algorithm (see subsection \ref{algorithmVEGA})
in the class \sourceRef{VegaAlgorithm}, in the method \sourceRef{mainPopulationToSegments},
individuals are divided into subpopulations according to the number of criteria. In each subpopulation,
individuals are assigned its designated criterion function,
or the evaluation is invalidated for those
who were evaluated by a different criterion in the previous generation.
In the method \sourceRef{segmentsToMainPopulation}, individuals
are evaluated within the subpopulation. The score is set for each as the ratio
of the calculated criterion value of the individual and the sum of all. Then the subpopulations
are merged into one.

\subsubsection{SPEA Algorithm}
The SPEA algorithm (see subsection \ref{algorithmSPEA}) in the class \sourceRef{SpeaAlgorithm}
is computationally the most complex due to the calculation of the non-dominated solution set and clustering.
The computation of strengths in the method named \sourceRef{computePowers}
first sets the strength of all individuals from the normal population to one.
Then it iterates through all individuals in the elite $E$ and in each step for each
individual $E_i$ it performs:
\begin{enumerate}
  \item records all individuals $D^i$ from the normal population that are dominated by individual $E_i$
  \item uses their count $|D^i|$ to determine the strength of individual $E_i$
  \item the strength of individual $E_i$ is added to all dominated individuals ${D^i}_j$
\end{enumerate}


For the implementation of clustering (described on page \pageref{secClustering}), the following
data types are used in the class \sourceRef{SpeaAlgorithm}:
\begin{verbatim}
  typedef std::vector<Sbox*> Cluster;
  class ClusterPair : public std::pair<Cluster*, Cluster*>;
  typedef std::vector<ClusterPair* > ClusterPairVector;
  typedef std::multimap<double, ClusterPair*> ClusterDistances;
  typedef std::map<Cluster*,  ClusterPairVector> PairsForCluster;
\end{verbatim}
A cluster is thus a list of pointers to individuals. The class \sourceRef{ClusterPair}
is auxiliary for calculating distances between clusters. It contains a flag
for memory management indicating whether it is still valid. The auxiliary type \sourceRef{ClusterPairVector} is used
to store both valid and invalid pairs for \sourceRef{ClusterDistances},
which maintains pairs sorted according to distances between clusters in the pair.
The double type key is used for continuous sorting, not for accessing elements using the double type.
And the auxiliary type \sourceRef{PairsForCluster} records all pairs in which a cluster
is contained. The entire clustering algorithm in the method \sourceRef{clustering} using these
types can be described:
\begin{enumerate}
  \item initialization of $pairsForCluster$ of type \sourceRef{PairsForCluster} \\ and
   $clusterDistances$ of type \sourceRef{ClusterDistances}
  \item first, the entire elite is traversed, and from each individual a cluster is created
   and added among the others
  \item traversing all $clusterDistances$ from the smallest distance, as long as the number
  of clusters is greater than the required elite size, invalid pairs are skipped
  \begin{enumerate}
    \item adding a new cluster containing both individuals from the current pair
    \item removing both old clusters
  \end{enumerate}
  \item generating a new elite from individuals closest to the cluster center
\end{enumerate}

When adding a cluster in the method \sourceRef{addCluster},
the distances of the new cluster with all existing clusters are calculated and
stored in $clusterDistances$.
Each new pair is also added to $pairsForCluster$ for both keys.

Removing a cluster in the method \sourceRef{removeCluster}:
When removing a cluster, pairs in which the cluster
is contained are invalidated and their list is removed for the current cluster in
$pairsForCluster$.

\subsubsection{Cartesian Genetic Programming, Random Search}
Cartesian genetic programming was divided into so-called "parallel random search"
and an acyclic gate graph, whose implementation is described below, because in this work it is
considered as an S-box representation. In the implementation of both random and parallel random search,
inheritance from the same base class as other evolutionary algorithms was used, for the purpose of easy comparison
thanks to the collection of statistics performed by GAlib. Therefore, they also use genome mutation
to generate further solutions.

\section{Implementation of S-box Representations}
For the implementation of S-box representations, specialized classes contained in GAlib were also used,
not only by deriving from \sourceRef{GAGenome}.
Figure \ref{sboxClassDiagram} shows the class diagram of representations and their derivation
from GAlib, which are marked in blue. This same color also marks the inheritance branch from
\sourceRef{Genome}.
In black are marked the custom classes of substitution box representations. This same color shows
the inheritance from the S-box base class. From the figure it is evident that multiple
inheritance was used. Originally, the class \sourceRef{Sbox} was designed as purely virtual,
but some functionalities that are common to all S-boxes were clearer to write into it than
somewhere outside the class. These methods include, for example, storing all possible inputs in an array,
an interface for calculating individual criteria, or performing multi-evaluation.
In contrast, the class \sourceRef{GAGenome} works only with one criterion.
In it, callbacks are defined as properties
for initialization, mutation, and crossover of the genome. In subclasses, these callbacks
are implemented and set.

\insertimage{sboxClassDiagram}{Class diagram of S-box representations with marked connection to GAlib}

While in \sourceRef{PermutationSboxTemplate} the output is the genes at corresponding positions, in
other representations it is necessary to perform a calculation. Therefore, in the base class there is the possibility to store these results
(and also criteria) and possibly invalidate them.
The code of \sourceRef{CgpGenome} was inspired by the implementation \cite{vasicek}.
The calculation of outputs for the i686 representation is given in \ref{secAbouti686}.

\section{Implemented Criteria and Their Evaluation}

In \sourceRef{multievaluate} (also in Fig. \ref{sboxClassDiagram}), the S-box outputs are first stored in an auxiliary dictionary
and from it the criterion values are calculated and then stored in the object so they do not have to be
calculated again.
In multi-criteria algorithms, it is therefore necessary to additionally call the method \sourceRef{invalidate} after mutation
of an individual so that the stored criterion values are invalidated.
Some calculated criteria need to be adjusted before being passed to the evolutionary algorithm; here are
their adjustments:
\begin{description}
  \item[bent score] the number of $w$ whose $|G(w)| = 1$
  \item[SAC order] one is added to the order because GAlib does not accept negative values
  \item[bent + SAC] if the bent score equals $2^n$, then the SAC order is added
  \item[differential probability] minimization performed using: $-\log_2(DP_{\max})$
  \item[linear probability] minimization performed using: $-\log_2(LP_{\max})$
  \item[bijective score] similar to bent score
\end{description}
Outputs are calculated using integer arithmetic. For GAlib, which works with
fitness in floating-point, conversion is done only at the end after calculating the entire criterion.
The outputs of some criteria, both during testing and final results, were compared with the mathematical
system SageMath and its functions for S-box\footnote{\url{http://www.sagemath.org/doc/reference/sage/crypto/mq/sbox.html}}.
Which was not suitable to use directly, because evolution for speed reasons runs in C++ and does not contain all
the criteria that the work deals with.

\section{Python Communication Interface}
The original prototype schema in Figure \ref{evolutionClassDiagram} was modified by
moving the classes for evolutionary algorithms to C++. Python therefore retains only the interface
for their superclass. This interface allows specifying in the constructor which algorithm
and genome should be used for searching.
When creating this interface, the ctypes library was used as much as possible, which allows
running functions in
the C language\footnote{C++ functions can be simply declared in an
\sourceRef{extern "C"} block}
located in a dynamically linked library. Primitive data types are
mostly converted to Python automatically. Sometimes it is necessary in Python to specify the type
of an argument or return value of a function, but this is not a problem if
the C functions do not change too much. The ctypes library
also supports C language structures.
Passing large data arrays is handled by
functions from the numpy library, which is used for processing results in experiments.
For passing strings, the Python C API was ultimately used.

In Appendix \ref{nutshell}, the interface is explained in a bit more detail from the perspective of its usage.

\section{Discussion of Parallel Computation Possibilities}
For parallelization to be worthwhile, code profiling was first performed,
during which parts of the code that took
the longest were found. What could no longer be accelerated further sequentially was selected for parallelization.
For an indicative estimate of whether parallelization was worthwhile, the
Unix tool \sourceRef{time} was used, which shows the consumed CPU
and real time of the program. CPU time for multiple processors (cores) is in
the output of this tool the sum from all used processors.
The CPU time of parallel execution and sequential execution had to be approximately the same.

Parallelization can be done at the level of evolution runs, with one evolution run running on each core/processor. Cases could also be handled where one run
finishes earlier and other runs (usually the last) could use the free time of an idle core, but this situation is rather exceptional for cases where
the program will be used. The runs share almost
no data with each other, and therefore parallelization on a multi-processor system is worthwhile. The only
thing they share and needed to be handled is the state of the pseudorandom generator.

For parallelization, OpenMP directives were used, which are more general,
more portable, and have a simpler and clearer notation
than, for example, POSIX threads. The code for POSIX threads is generated there by the GCC compiler, other
compilers can generate different code specific to a particular processor architecture.

Thanks to OpenMP, it was possible to mark as
private\footnote{private from the perspective of sharing by multiple processors}
variables that store
the state of the pseudorandom generator in GAlib:
\begin{verbatim}
  #pragma omp threadprivate(seed, iseed, idum2, iy, iv, idum)
\end{verbatim}
This modification did not interfere with
the GAlib interface and it is possible when compiling for sequential execution
to use the original GAlib.

Another possible parallelization would be to calculate the objective
functions of individuals within the population in parallel. The class \sourceRef{GAPopulation} maintaining the list of individuals
in the population would have to be implemented safely for
use by multiple threads, i.e., with added synchronization.
\pagebreak
By the considered parallelization, the following code is meant:
\begin{verbatim}
// definition of population evaluator
void OmpPopulationEvaluator(GAPopulation & p) {
  const int size = p.size();
  #pragma omp parallel for shared(p) schedule(guided)
  for(int i=0; i< size; i++)
    p.individual(i).evaluate();
}

// setting the evaluator e.g. in the factory method of GAGeneticAlgorithm object
GAPopulation initPopulation(ga.population());
initPopulation.evaluator(OmpPopulationEvaluator);
ga.population(initPopulation);
\end{verbatim}

Even without all necessary synchronizations, this parallel solution ultimately still took longer CPU time
than the sequential solution, so not much would be gained by this parallelization for most criteria,
unless the calculation of some criterion were more complex, e.g., $LP_{\max}$.
